\section{Challenge Format}
\label{sec:challenge_format}

The PZ data challenge comprises a series of sets of tasks for participants.
Submissions will be evaluated to determine how ready various algorithms are
to be used for cutting-edge analysis based on how well they perform on
the various tasks. Readiness will be evaluated on a few different
fronts: 1) Does the algorithm meet performance requirements? 2) Is
it robust, flexible, and relatively easy to use on different datasets?
3) Is it scalable up to the scales we will need to use it at?

This document and the associated web pages describe the data being
provided to participants, the tasks they will be asked to perform, the
expected format for submission and the metrics by which the algorithm
readiness will be evaluated.


\subsection{Scope and Timeline}

The data challenge will include two major parts, with a set of
tasks emulating increasingly realistic scenarios in each part. The
first part, $p(z)$ estimation, will focus on estimating the redshift of
individual objects. The second part, tomography and $n(z)$ estimation,
will focus on assigning objects to tomographic bins and estimating the
distribution of redshifts in each bin.

The data challenge will run from April 13, 2026 to September 2026.
The first set of data and tasks related to $p(z)$ estimation will be
released on April 13 and will close on July 17, 2026.   A second set
of data and tasks related to more realistic $p(z)$ estimation
scenarios and $n(z)$ estimation will be released on June 1 and will
close in September 2026.

Preliminary results will be released in August 2026, with a
technical note summarizing those results to follow shortly thereafter and
a comprehensive journal publication to follow later.


\subsection{Installing and setting up the \texttt{pz\_data\_challenge}
  package}

The \texttt{pz\_data\_challenge} package will provide participants with tools
to access data, set up submissions, estimate performance metrics and
format submissions.  This can be set up with a few small
variants on the standard \texttt{GitHub} package setup procedure.
Before starting you should pick a name for your submission, e.g., ``example''.

\begin{verbatim}
# Create a conda environment
conda create --name pzdc python=3.13

# Clone the pz_data_challenge repository (or your fork of the repository)
git clone git@github.com:LSSTDESC/pz_data_challenge.git
# or git clone https://github.com/LSSTDESC/pz_data_challenge.git

# Go into the directory
cd pz_data_challenge

# Install the code in "editable" mode
pip install -e ".[dev]"

# Use the provided script to set up your submission.
# Here you should provide the name of your submission
python scripts/prepare_submission.py <submission_name>
\end{verbatim}

This final step will copy the input data files to
\texttt{pz\_data\_challenge/public}, and set up the three files you will
need to submit your entry.

The notebooks in the \texttt{pz\_data\_challenge/nb} area give examples
of how to access the data and create some of the diagnositic plots
that were used to validate the data.


\subsection{Submission mechanism}

Submission will take the form of pull request in the
\texttt{pz\_data\_challenge} repository.  Detailed instructions on how
to submit an entry are provided in
Sec.~\ref{sec:challenge_submissions} of this document.

